%% ---------------------------------------------------------------------
\section{Sistemas interconectados}
\label{sec:av-inter}
%% ---------------------------------------------------------------------

En esta sección $x_1\in\R^{n_1}$ y $x_2\in\R^{n_2}$ son los estados de \emph{dos subsistemas}; en los ejemplos son escalares, pero el razonamiento es por subsistemas. Para la cascada
\[
  \dot x_1=f_1(x_1,x_2),\qquad\dot x_2=f_2(x_2),
\]
con el origen de $\dot x_2=f_2(x_2)$ GAS y el de $\dot x_1=f_1(x_1,0)$ GAS, la intuición dice: ``como $x_2\to0$, a la larga $x_1$ ve su sistema sin entrada''. El problema está en ``a la larga'': antes de que $x_2$ se apague, $x_1$ puede escapar en tiempo finito, y entonces no hay ``larga''.

\begin{ejemplo}[label=ej:av-cascada]{GAS más GAS no da GAS \citep[Ej.~4]{Arc18}}
Estudie $\dot x_1=-x_1+x_1^2x_2$, $\dot x_2=-x_2$.
\solucion Con $x_2=x_{20}e^{-t}$ y $z=1/x_1$, $\dot z=z-x_{20}e^{-t}$, lineal, con solución $z=(z_0-x_{20}/2)e^t+(x_{20}/2)e^{-t}$. $x_1$ escapa cuando $z$ se anula, lo que ocurre en
\[
  t_e=\frac12\ln\frac{x_{10}x_{20}}{x_{10}x_{20}-2}
\]
si y solo si $x_{10}x_{20}>2$. Desde $(2,2)$, $t_e=\tfrac12\ln2=0.347$ (\cref{fig:av-cascada}).
\end{ejemplo}

\begin{trampa}[label=tr:av-cascadaerr]{``$\dot x_1=-x_1+x_1x_2^2$ también escapa''}
\plausible{Si $x_1^2x_2$ produce escape, $x_1x_2^2$ debería producirlo también; así lo decían Handout~7 y las diapositivas.}
Es falso. La ecuación de $x_1$ es lineal en $x_1$ y
\[
  x_1(t)=x_{10}\exp\Bigl(-t+\frac{x_{20}^2}{2}\bigl(1-e^{-2t}\bigr)\Bigr)\to0 ,
\]
con $\abs{x_1(t)}\le\abs{x_{10}}e^{x_{20}^2/2}$: el origen es GAS. El transitorio puede ser grande (desde $(1,3)$, $x_1$ sube hasta $18.2$), pero no escapa.
\correcto{Sin embargo el subsistema $x_1$ no es ISS respecto a $x_2$: con $x_2\equiv c>1$, $\dot x_1=(c^2-1)x_1$ crece. ISS de $x_1$ es \emph{suficiente} para que la cascada sea GAS, no necesaria.}
\end{trampa}

\begin{proposicion}[label=prop:av-cascadaiss]{Cascadas ISS \citep{Son89a}; \citep[Teorema~7.9]{KeBr23}}
Si en $\dot x_1=f_1(x_1,x_2)$, $\dot x_2=f_2(x_2,w)$ el subsistema $x_1$ es ISS con $x_2$ como entrada y el subsistema $x_2$ es ISS con $w$ como entrada, la cascada es ISS con entrada $w$.
\end{proposicion}

\begin{corolario}[label=cor:av-cascadagas]{\citep[Lema~4.7]{Kha02}}
Si el origen de $\dot x_2=f_2(x_2)$ es GUAS y $\dot x_1=f_1(x_1,x_2)$ es ISS respecto a $x_2$, el origen de la cascada es GUAS.
\end{corolario}

\begin{demostracion}
Por ISS de $x_1$, para todo $s\in[t_0,t]$,
$\norm{x_1(t)}\le\beta_1(\norm{x_1(s)},t-s)+\gamma_1\bigl(\sup_{s\le\tau\le t}\norm{x_2(\tau)}\bigr)$, y $\norm{x_2(\tau)}\le\beta_2(\norm{x_2(t_0)},\tau-t_0)$. Con $s=t_0$ sale el acotamiento (y la estabilidad): $\norm{x_1(t)}\le\beta_1(\norm{x_1(t_0)},0)+\gamma_1(\beta_2(\norm{x_2(t_0)},0))$. Con $s=(t+t_0)/2$, \emph{el punto medio},
\[
  \norm{x_1(t)}\le\beta_1\Bigl(\norm{x_1(s)},\tfrac{t-t_0}2\Bigr)+\gamma_1\Bigl(\beta_2\bigl(\norm{x_2(t_0)},\tfrac{t-t_0}2\bigr)\Bigr)\to0 :
\]
la mitad del tiempo para que $x_2$ se apague y la otra mitad para que $x_1$ olvide su condición en $s$.
\end{demostracion}

\begin{figure}[ht]
\centering
\includegraphics[width=\linewidth]{cap10/f_cascada}
\caption{Izquierda: la cascada del \cref{ej:av-cascada} desde $(2,2)$ escapa en $t=0.347$. Centro: $\dot x_1=-x_1+x_1x_2^2$ desde $(1,3)$ sube hasta $18.2$ y vuelve a cero. Derecha: la cascada ISS del \cref{ej:av-cascadaiss} con $w=0.3\sin t$, acotada.}
\label{fig:av-cascada}
\end{figure}

\begin{ejemplo}[label=ej:av-cascadaiss]{$\dot x_1=-x_1+x_2^2$, $\dot x_2=-x_2+w$}
Pruebe que es ISS de dos maneras: por cascada y con una función de Lyapunov compuesta.
\solucion \emph{Por cascada.} Con $V=\tfrac12x_1^2$, $\Vdot=-x_1^2+x_1x_2^2\le-(1-\theta)x_1^2$ si $\abs{x_1}\ge x_2^2/\theta$: el subsistema $x_1$ es ISS respecto a $x_2$ con $\gamma_1(r)=r^2/\theta$; el $x_2$ lo es respecto a $w$. Por la \cref{prop:av-cascadaiss}, la cascada es ISS.
\emph{Con $V$ compuesta.} Se prueba $V=\tfrac12x_1^2+\tfrac a4x_2^4$ (cuártica en $x_2$ porque el acoplamiento $x_1x_2^2$ es de orden~4 al completar el cuadrado). Con $w=0$,
\[
  \Vdot=-\Bigl(x_1-\frac{x_2^2}2\Bigr)^2-\Bigl(a-\frac14\Bigr)x_2^4 ,
\]
negativa definida si $a>\tfrac14$. Con $a=2$ y $w$, $\Vdot=-x_1^2+x_1x_2^2-2x_2^4+2x_2^3w$, y dos usos de Young \eqref{eq:av-young}: $x_1x_2^2\le\tfrac12x_1^2+\tfrac12x_2^4$ y, con $p=\tfrac43$, $q=4$, $\lambda^4=8/27$, $\abs{x_2}^3\abs w\le\tfrac12x_2^4+\tfrac{27}{32}w^4$. Queda
\[
  \Vdot\le-\tfrac12x_1^2-\tfrac12x_2^4+\tfrac{27}{16}w^4 ,
\]
la forma disipativa: ISS. (Con $a=1$, Handout~7 tenía un signo duplicado; lo correcto es $\Vdot=-\tfrac12(x_1-x_2^2)^2-\tfrac12(x_1^2+x_2^4)+x_2^3w$.)
\end{ejemplo}

\codigorepo{cap10_avanzada/C10_Ej14_Cascadas.m}{cascadas}{Escape, transitorio grande y cascada ISS}

\subsection{Retroalimentación y pequeña ganancia}

Si en lugar de una cascada hay una \emph{retroalimentación}, $\dot x_1=f_1(x_1,x_2)$, $\dot x_2=f_2(x_2,x_1)$, con cada subsistema ISS respecto al otro (ganancias $\gamma_1$, $\gamma_2$), la ISS de cada uno ya no basta: la señal da vueltas por el lazo y puede amplificarse.

\begin{teorema}[label=thm:av-pequenaganancia]{Pequeña ganancia \citep{JTP94}}
Si $\gamma_1\circ\gamma_2(r)<r$ para todo $r>0$, el origen de la interconexión es GAS.
\end{teorema}

La parte fácil de la prueba muestra la idea. Si las soluciones están acotadas y $L_i=\limsup\norm{x_i(t)}$, reiniciando el reloj como en la \cref{prop:av-issconsec}, $L_1\le\gamma_1(L_2)$ y $L_2\le\gamma_2(L_1)$, así que $L_1\le\gamma_1(\gamma_2(L_1))<L_1$ si $L_1>0$: imposible. \alerta{Una vuelta completa por el lazo contrae.} Lo difícil es probar el acotamiento y la estabilidad.

\begin{ejemplo}[label=ej:av-pequenaganancia]{$\dot x_1=-x_1+a\,x_2$, $\dot x_2=-x_2+b\sin x_1$}
¿Para qué $a,b$ es GAS el origen?
\solucion Por variación de constantes, $\gamma_1(r)=\abs ar$, y como $\abs{\sin x_1}\le\abs{x_1}$, $\gamma_2(r)=\abs br$. La condición es $\abs{ab}<1$: GAS.
\emph{Con $ab>1$.} Los equilibrios cumplen $x_1=ab\sin x_1$, con soluciones $\pm\eq_1\neq0$; con $a=1$, $b=1.2$, $\eq_1=1.027$. El jacobiano tiene valores propios $-1\pm\sqrt{ab\cos x_1}$: en el origen, $-1\pm\sqrt{1.2}$, silla; en $\pm\eq_1$, $ab\cos\eq_1=0.62<1$, nodos estables. El origen deja de ser GAS.
\emph{Con $ab=-3$: GAS aunque $\abs{ab}>1$.} $x_1=-3\sin x_1$ solo tiene la solución $0$; las soluciones son últimamente acotadas; la divergencia es $-2$, así que no hay órbitas cerradas (Bendixson, \cref{cap:orbitas}); el único equilibrio es un foco estable ($-1\pm j\sqrt3$), y por Poincaré--Bendixson es el $\omega$-límite de todas las trayectorias. \alerta{La pequeña ganancia es suficiente, no necesaria.}
\end{ejemplo}

\begin{figure}[ht]
\centering
\includegraphics[width=\linewidth]{cap10/f_pequena_ganancia}
\caption{El \cref{ej:av-pequenaganancia}. Izquierda: $ab=0.8$, todas las trayectorias van al origen. Centro: $ab=1.2$, el origen es una silla y aparecen dos nodos estables en $x_1=\pm1.027$. Derecha: $ab=-3$, viola la pequeña ganancia y aun así el origen es GAS (foco).}
\label{fig:av-pequenaganancia}
\end{figure}

\codigorepo{cap10_avanzada/C10_Ej15_PequenaGanancia.m}{ganancia}{Los tres retratos de fase de la retroalimentación}

\begin{pausa}
$\dot x_1=-x_1+2x_2$, $\dot x_2=-x_2+x_1$: ¿cada subsistema es ISS respecto al otro? ¿Se cumple la pequeña ganancia? ¿Es GAS el origen?
\respuesta{Cada subsistema es lineal estable y, por \eqref{eq:av-isslineal}, ISS con ganancias $2$ y $1$. Pero $\gamma_1\gamma_2=2>1$: la pequeña ganancia falla. Y el origen no es GAS: la matriz $\bigl[\begin{smallmatrix}-1&2\\1&-1\end{smallmatrix}\bigr]$ tiene el valor propio $-1+\sqrt2>0$. Aquí la falla de la condición sí se traduce en inestabilidad.}
\end{pausa}
